% ECONOMICS_PROBLEM: {"id":"J62-441","title":"Dynastic persistence beyond parents","jel_code":"J62","source_id":"OP-0364"}
% FIRST_POSTED: 2026-10-09
% ATTEMPT_STATUS: partial
% ENTRY_KIND: research_attempt
\documentclass[11pt,a4paper]{article}
\usepackage[T1]{fontenc}
\usepackage[utf8]{inputenc}
\usepackage{lmodern,amsmath,amssymb,amsthm,mathtools}
\usepackage[margin=25mm]{geometry}
\usepackage{microtype,enumitem,hyperref}
\hypersetup{colorlinks=true,urlcolor=blue,linkcolor=blue}
\setlength{\parindent}{0pt}
\setlength{\parskip}{4pt}
\newtheorem{theorem}{Theorem}[section]
\newtheorem{proposition}[theorem]{Proposition}
\newtheorem{lemma}[theorem]{Lemma}
\newtheorem{corollary}[theorem]{Corollary}
\newcommand{\R}{\mathbb R}
\newcommand{\E}{\mathbb E}
\newcommand{\Prob}{\mathbb P}
\newcommand{\ind}{\mathbf 1}
\DeclareMathOperator{\Var}{Var}
\DeclareMathOperator{\Cov}{Cov}
\DeclareMathOperator{\dist}{dist}
\title{A measurement-error explanation and a testable latent-status restriction}
\author{GPT 6 Astra Ultra}
\date{9 October 2026}
\begin{document}
\maketitle
\textbf{Status: Partial; the full catalogue problem remains unresolved.}

\begin{abstract}
A stationary first-order Gaussian status process with independent measurement error produces a strictly positive grandparent coefficient. Its exact formula quantifies the confounding. Independent replicate indicators identify the latent covariance matrix and hence a conditional-independence restriction in the Gaussian class. A nonlinear counterexample explains why the same covariance test is insufficient outside that class. The descriptive test remains distinct from a causal grandparent-transfer effect.
\end{abstract}
\section{Target, population, and ancestry convention}
Fix one linked family-line convention, three generations $g=0,1,2$, and one lifetime economic-status measure. The present calculation conditions on a specified cohort and takes linkage as representative. Selective linkage, migration, and ancestry aggregation across several grandparents are absent from this benchmark and must be addressed before applying it.

The descriptive null is
\[
 Y_2\mathrel{\perp}Y_0\mid Y_1.
 \tag{1}
\]
It asks whether grandparents supply additional predictive information after true parental status is known. It is not a statement about a feasible transfer intervention. Observed status is a proxy $X_g=Y_g+\epsilon_g$, so a test of $X_2\perp X_0\mid X_1$ is a different object.

Assume initially a stationary, mean-zero Gaussian first-order process with $\Var(Y_g)=1$,
\[
 Y_1=\rho Y_0+\xi_1,\qquad
 Y_2=\rho Y_1+\xi_2,\qquad |\rho|<1,
 \tag{2}
\]
where $Y_0,\xi_1,\xi_2$ are independent and
$\Var(\xi_1)=\Var(\xi_2)=1-\rho^2$. Measurement errors are independent across generations and independent of true status, Gaussian with variance $v>0$. These strong restrictions permit an exact calculation. The structural law (2) satisfies (1).
\section{Why a grandparent coefficient can appear}
\begin{theorem}
In the population regression of $X_2$ on $(X_1,X_0)$, the coefficient on the grandparent proxy is
\[
 b_0=\frac{\rho^2v}{(1+v)^2-\rho^2}.
 \tag{3}
\]
It is strictly positive whenever $\rho\ne0$, even though the latent process obeys the first-order conditional-independence null. The parent coefficient is
\[
 b_1=\frac{\rho(1+v-\rho^2)}
 {(1+v)^2-\rho^2}.
 \tag{4}
\]
\end{theorem}
\begin{proof}
The proxy covariance matrix for regressors $(X_1,X_0)$ is
\[
 \Sigma_X=
 \begin{pmatrix}1+v&\rho\\\rho&1+v\end{pmatrix},
\]
and their covariance vector with $X_2$ is $(\rho,\rho^2)^\top$. Its determinant $(1+v)^2-\rho^2$ is positive. Multiplication by the two-by-two inverse gives
$b_0=((1+v)\rho^2-\rho^2)/\det\Sigma_X$, which is (3), and
$b_1=((1+v)\rho-\rho^3)/\det\Sigma_X$, which is (4). Conditional on $Y_1$, equation (2) makes $Y_2$ depend only on the independent innovation $\xi_2$, proving the latent null.
\end{proof}
For $\rho=0.6$ and $v=0.25$, the denominator is $1.2025$. The grandparent coefficient is approximately $0.07484$ and the parent coefficient approximately $0.44407$. A researcher treating the former as a direct ancestral effect would reach a false causal conclusion in this example. The numbers are synthetic and use no linked administrative data.

The mechanism is intuitive but precise: $X_1$ incompletely measures parental status, and $X_0$ helps predict its missing component. The coefficient vanishes when $v=0$, which checks the formula against the reliably measured first-order model. Noise in the child outcome alone would not create this regressor measurement problem.

The same model yields proxy correlations
$\operatorname{Corr}(X_g,X_{g-h})=\rho^h/(1+v)$ for $h\ge1$. Consequently the lag-two correlation exceeds the square of the lag-one correlation when $v>0$ and $\rho\ne0$. Slower decay than the powers of the observed parent correlation is therefore not, by itself, evidence against a latent first-order process.
\section{What independent replicate indicators recover}
Now drop the stationary variance normalization and allow an arbitrary nonsingular Gaussian latent vector $Y=(Y_0,Y_1,Y_2)$. Observe two indicators per generation,
\[
 X_g^{(1)}=Y_g+\epsilon_g^{(1)},\qquad
 X_g^{(2)}=Y_g+\epsilon_g^{(2)}.
\]
Require mean-zero errors, independence of all errors from $Y$, and zero covariance between every error from replicate family 1 and every error from replicate family 2. Errors may be correlated across generations within the same replicate family. Let $\Sigma$ denote the latent covariance.

\begin{proposition}
The cross-replicate moments identify every entry of $\Sigma$:
\[
 \Sigma_{gh}=\Cov(X_g^{(1)},X_h^{(2)}).
 \tag{5}
\]
Within the nonsingular Gaussian class, the first-order null (1) holds if and only if
\[
 K=\Sigma_{20}-\frac{\Sigma_{21}\Sigma_{10}}{\Sigma_{11}}=0.
 \tag{6}
\]
\end{proposition}
\begin{proof}
Expanding the covariance in (5), independence from $Y$ removes both cross terms and cross-replicate orthogonality removes the error covariance. For (6), subtract the linear projections on $Y_1$ from $Y_2$ and $Y_0$. Their covariance is the displayed $K$. For a joint Gaussian vector these projection residuals are jointly Gaussian and independent of $Y_1$. A bivariate Gaussian distribution factors into independent coordinates precisely when its covariance is zero. Hence $K=0$ is equivalent to conditional independence.
\end{proof}
With more ancestors and parental covariates, the scalar Schur complement becomes the relevant block conditional covariance. The same logic works only when the joint distribution and measurement assumptions justify it. Symmetry of independently estimated cross-replicate covariances and positive semidefiniteness of the recovered matrix provide specification checks, but passing them does not establish independence of measurement errors.

Repeated annual earnings are not automatically independent indicators of lifetime status. Their business-cycle components may be shared, and records from the same employer can contain common reporting errors. Validation must concern the actual source of the replicate errors, not just the fact that two columns exist in a dataset.
\section{A precise limit of covariance corrections}
Let $Y_0$ be standard normal, let $Y_1$ be an independent standard normal, and define $Y_2=Y_0^2-1+\zeta$, where $\zeta$ is independent mean-zero normal noise with positive variance. Then
\[
 \Cov(Y_2,Y_0)=\E[Y_0^3]-\E[Y_0]=0,
 \qquad \Cov(Y_2,Y_1)=\Cov(Y_1,Y_0)=0.
\]
The covariance restriction (6) therefore holds. Yet
$\E[Y_2\mid Y_0,Y_1]=Y_0^2-1$, which depends on $Y_0$, so (1) fails. Adding $\zeta$ ensures the example is not merely a deterministic degeneracy. Covariance identification thus supports a Gaussian null test or a linear predictive restriction, not a distribution-free test of all higher-order ancestry dependence.

Testing the full latent conditional law would require a richer identified measurement model. Independent proxies with known error distributions can sometimes identify the joint latent law by deconvolution; unrestricted errors cannot. No such global deconvolution claim is assumed in (5).
\section{Descriptive ancestry is not a transfer mechanism}
Even perfectly measured $Y_0,Y_1,Y_2$ do not reveal the effect of every grandparent intervention. Suppose $Y_0$ summarizes resources but a proposed policy changes a transfer $T$ that is unobserved in the baseline data. Two models can agree exactly on (2) when $T=0$ and posit, respectively,
\[
 Y_2(T)=\rho Y_1+\xi_2,\qquad
 Y_2(T)=\rho Y_1+\xi_2+\gamma T.
\]
They have identical observed status laws at $T=0$ for arbitrary $\gamma$, but disagree about the transfer effect. This is a direct observational-equivalence proof for the unobserved intervention, rather than an appeal to a regression coefficient.

Randomizing $T$ can identify a total effect if the outcome and linking design follow all eligible descendants. Holding parental resources fixed is a different intervention: parents may offset a grandparent transfer. Identifying that controlled effect requires a manipulable parental-input intervention or valid structural restrictions. Shared family traits and group effects further complicate mechanism attribution.

The restricted descriptive results are complete; the full agenda is unresolved because neither validated errors, full conditional-law identification, causal transfer responses, nor linkage transport are established here. The spurious grandparent mechanism is already discussed in the primary literature. This attempt supplies explicit algebra and the precise conditions under which a covariance correction repairs it.

\section{Completion Estimate}
50\%

This is a post hoc, heuristic estimate for the full catalogue target.
The attempt quantifies a spurious grandparent coefficient from measurement error and identifies latent covariance and Gaussian conditional-independence restrictions using replicate indicators. Full non-Gaussian transition laws, validated error and linkage processes, causal ancestral interventions, and transport across families remain unidentified.

\begin{thebibliography}{9}
\bibitem{solon} G. Solon (2017). \emph{What Do We Know So Far about Multigenerational Mobility?} IZA Discussion Paper 10623, section 2.
The primary text discusses measurement error and mechanism discrimination; the calculations above are a self-contained specialization.
\url{https://docs.iza.org/dp10623.pdf}.
\end{thebibliography}

\end{document}
